%% Document Layout and Structure %%
\documentclass[12pt,a4paper]{report}
\usepackage[T1]{fontenc} 
\usepackage[ngerman]{babel} % comment out if thesis written in English
\usepackage[left=2.5cm,right=2.5cm,top=2.5cm,bottom=2.0cm,bindingoffset=8mm,includehead,includefoot]{geometry} % check with the printing shop if more space for binding (bindingoffset) is needed

\usepackage[onehalfspacing]{setspace} % spacing between lines
\parindent 0pt % no indentation after paragraph

\usepackage{titlesec} % for customization of chapters, sections and subsections
\titleformat{\chapter}[display]
  {\bfseries\filright}
  {\hrule height 0pt \huge\thechapter\hrule}
  {-\baselineskip}
  {\Huge\strut}
\titlespacing*{\chapter}
  {0pt}
  {-\topskip}
  {2\normalbaselineskip}

%% Colours %%
\usepackage{xcolor} % colour management and manipulation, supports a large number of colour models

%% Code Snippets %%
\usepackage{listings} % useful for displaying code snippets etc.

%% Hyperlinks %%
\usepackage[hidelinks]{hyperref} % hidelinks removes borders around the hyperlinks

%% Graphics %%
\usepackage{graphicx} % to include graphics or images
\usepackage{svg}

%% Lorem Ipsum %% 
\usepackage{lipsum} % to have filler text for this template

%% Acronyms %%
\usepackage[printonlyused,withpage]{acronym} % both options optional; for more features (e.g. acronyms & glossary) you can instead use the package glossaries (see e.g. https://de.overleaf.com/learn/latex/Glossaries)

%% Bibliography & Citation %%
% For the package biblatex a bibliography data file is needed; use \printbibliography at the end of the document to display the references cited in your thesis (thesis.tex)
\usepackage[style=authoryear,citestyle=authoryear-icomp,uniquename=minfull,maxbibnames=99,maxcitenames=2]{biblatex} % you can change the options (e.g. style & citestyle)
\addbibresource{bibliography.bib} % imports the bibliography data file

\usepackage{csquotes} % for quotations